\sect{Конвейер обработки в реальном времени}{realtime}
%
Архитектура конвейера определяется требованием выдавать
аккомпанемент с задержкой не более $20$\,мс относительно момента
звукоизвлечения солиста. Этот бюджет распределяется между четырьмя
блоками, изображёнными на \Figref{pipeline}: входной модуль,
блок очистки, гибридный трекер HMM/OLTW и потоковый рендерер.

\par\addvspace{2pt}
{\centering
\refstepcounter{figure}\figlab{pipeline}%
\resizebox{\columnwidth}{!}{% =====================================================================
% TikZ figure: end-to-end real-time pipeline (BLACK & WHITE only)
% Compact 3+3 layout to fit a single REVTeX two-column column.
% Used in: sections/05-realtime-pipeline.tex
% =====================================================================
\begin{tikzpicture}[
    node distance = 4mm and 4mm,
    every node/.append style={font=\scriptsize},
    block/.style ={rectangle, draw=black, thick, rounded corners=2pt,
                   minimum height=8mm, minimum width=18mm,
                   align=center, font=\scriptsize},
    arrow/.style ={-Stealth, thick}]

  \node[block]                 (in)    {MIDI / audio\\input};
  \node[block, right=of in]    (clean) {Cleaning\\\& buffer};
  \node[block, right=of clean] (hmm)   {HMM\\tracker};

  \node[block, below=8mm of hmm]   (oltw) {OLTW\\fall-back};
  \node[block, left=of oltw]       (rl)   {RL tempo\\anticip.};
  \node[block, left=of rl]         (out)  {Orchestra\\renderer};

  \draw[arrow] (in)    -- (clean);
  \draw[arrow] (clean) -- (hmm);
  \draw[arrow] (hmm)   -- (oltw);
  \draw[arrow] (oltw)  -- (rl);
  \draw[arrow] (rl)    -- (out);

  \draw[arrow, dashed] (hmm.south west) to[bend right=15]
    node[left, font=\tiny, pos=0.55]
        {$\max\alpha_t<\theta_{\text{conf}}$} (oltw.north west);

\end{tikzpicture}
}\par
\vspace{2pt}
\footnotesize
\figurename~\thefigure. Архитектура конвейера обработки в реальном
времени. Сплошные стрелки~--- основной поток данных, штриховая~---
переключение на резервный трекер при низкой уверенности HMM.\par}
\addvspace{4pt}

Входной модуль работает в двух режимах. В режиме MIDI используется
библиотека \texttt{python-rtmidi}, доставляющая события NoteOn и
NoteOff с точностью до миллисекунды. В аудиорежиме применяется
PortAudio с буфером $128$ сэмплов при частоте $44{,}1$\,кГц
(теоретическая задержка $2{,}9$\,мс плюс время прохождения через
аудиодрайвер ОС, которое в Windows~ASIO составляет $4$--$6$\,мс).
Принятые события маркируются временной меткой источника и помещаются
в кольцевой буфер с потокобезопасным доступом.

Очистка входных данных выполняет три функции. Во-первых, отсеиваются
события с длительностью менее $25$\,мс (чаще всего следствие случайных
касаний клавиш); во-вторых, восстанавливаются октавные ошибки путём
сравнения с ожидаемой высотой текущего HMM-состояния (если кандидат
$q_j+12$ или $q_j-12$ ближе к $p_{\hat\varphi}$, чем сам $q_j$,
кандидат заменяется); в-третьих, для аудиорежима производится
агрегация смежных onset-кандидатов в окне $40$\,мс. Подобная обработка
снижает количество ложных переходов в HMM в среднем на $11$\,\% по
сравнению с подачей сырых событий.

Гибридный трекер реализует схему из~\Secref{background}. На каждом
такте выполняется обновление forward-распределения по
\Eqref{forward-recursion}, после чего вычисляется уверенность
$c_t=\max_i\alpha_t(i)$. Если $c_t\ge\theta_{\text{conf}}$, выходом
служит позиция $\hat\varphi(t)$, выданная HMM. В противном случае
управление переходит к OLTW, который инициализируется текущим
наилучшим кандидатом HMM, и одновременно запускается процедура
повторной синхронизации: HMM переинициализирует $\alpha$ в
сглаженной окрестности $\hat\varphi$ с шириной $\pm 4$ ноты, что
позволяет восстановиться после кратковременных ошибок исполнителя.

Выходной модуль отвечает за фактическое воспроизведение оркестровой
партии. Используется собственный сэмплер на основе библиотеки
FluidSynth, поддерживающий растяжение времени без изменения тембра
методом фазового вокодера. Растяжение применяется на коротких блоках
($512$ сэмплов), что эквивалентно временно́му разрешению $11{,}6$\,мс.
Темп блока задаётся не последним наблюдением, а \textit{предсказанием}
RL-модуля (см.~\Secref{rl}) на горизонте $200$--$500$\,мс, что
позволяет компенсировать суммарную задержку конвейера и достичь
эффективной задержки выходного сигнала менее $20$\,мс.

\par\addvspace{2pt}
{\centering
\refstepcounter{table}\tablab{latency-budget}%
\footnotesize
\begin{tabular}{@{}lc@{}}
  \toprule
  Этап обработки                              & Бюджет, мс \\
  \midrule
  Захват события (MIDI/audio)                 & 1--6 \\
  Очистка и буферизация                       & 1--2 \\
  Обновление HMM (Forward + softmax)          & 2--3 \\
  Резервный OLTW (включается по уверенности)  & 3--5 \\
  Запрос RL-предсказания темпа                & 1--2 \\
  Растяжение и синтез блока FluidSynth        & 6--8 \\
  \midrule
  Суммарная end-to-end задержка               & 14--26 \\
  \bottomrule
\end{tabular}\par
\vspace{2pt}
\tablename~\thetable. Распределение задержек по этапам конвейера.
Целевое значение end-to-end задержки~--- не более $20$\,мс.\par}
\addvspace{4pt}

Контроль соблюдения бюджета осуществляется встроенным профайлером,
который раз в секунду пишет в лог время выполнения каждого блока.
Превышение бюджета на любом из этапов помечается как событие
деградации и используется при анализе экспериментов
(\Secref{experiments}).
